\section{Random routing and independent tests}
\label{sec:routing}

Keep the tree, grids, and windows of Section~\ref{sec:windows}, and fix
\(p\in(0,1)\). We construct a random periodic set of expected density
\(p\), routing each point through the tree with one child at each internal vertex
designated as the default. When the route of a stable center makes a
default choice, the first such vertex supplies many tests for membership
of nearby points. We compute the probability of having no default choice
and show that, when the tests are available, their conditional failure
probability at each fixed scale decays exponentially in the window length.
The dependence on \(q\) enters through the grids and separation estimates
already established.

\subsection{The random tables}

For each edge \(e=(P,i)\) with \(i<M\), let
\[
 S_e:\{0,\ldots,N_{b_e}-1\}\longrightarrow\{0,1\}
\]
be a table of independent fair bits. Child \(M\) is the default child
and has no selector table. For each leaf \(L\), let \(b(L)\) be the
window endpoint of its incoming edge, and let
\[
 T_L:\{0,\ldots,N_{b(L)}-1\}\longrightarrow\{0,1\}
\]
be a table of independent Bernoulli-\(p\) bits. All entries of all
selector tables \(S_e\) and terminal tables \(T_L\) are mutually
independent. There are finitely many entries, so they define a finite
product probability space \(\Omega\), with probability \(\Prob\) and
expectation \(\E\). Write \(\mathcal A\) for the sigma-field generated
by all entries of all selector tables.

Given an outcome and \(z\in\R\), route \(z\) from the root as follows.
At an internal vertex \(P\), choose the first child \(P_i\), \(i<M\),
whose selector satisfies
\(S_{(P,i)}(J_{b_{(P,i)}}(z))=1\). If all these values are zero, choose
\(P_M\). Let \(L(z)\) denote the leaf reached and define
\begin{equation}
 \label{eq:random-set}
 B=\bigl\{z\in\R:
 T_{L(z)}(J_{b(L(z))}(z))=1\bigr\}.
\end{equation}
Every key is periodic, and every grid is refined by the finest grid in
the construction. Thus each realization of \(B\) is \(1\)-periodic
and is a union of finitely many half-open cells per period.

\noindent\textit{Remark (notation).} For each fixed $z\in\R$, $L(z)$ is the selector-dependent leaf reached from the root by the stated routing rule; it is $\mathcal A$-measurable and is fixed upon conditioning on $\mathcal A$. $B=\{z\in\R:T_{L(z)}(J_{b(L(z))}(z))=1\}$, with routing and terminal membership evaluated in the same table outcome, and each realization is $1$-periodic and a union of finitely many half-open grid cells per period. In Proposition~\ref{prop:periodic}: After the expectation estimates, one table outcome is fixed with $\rho(B^+)+\rho(R)\le5p$, and the subsequent construction uses its realization of $B$.


\begin{lemma}[Mean density]
 \label{lem:density}
The random set \(B\) in \eqref{eq:random-set} satisfies
\(\E\rho(B)=p\).
\end{lemma}

\begin{proof}
For fixed \(z\), conditioning on \(\mathcal A\) fixes
its leaf and terminal address. The terminal entry at that address still
has Bernoulli-\(p\) law, so \(\Prob(z\in B\mid\mathcal A)=p\).
Integration over one period, interchanged with the finite weighted sum
over \(\Omega\), gives
\[
 \E\rho(B)=\int_0^1\Prob(z\in B)\,dz=p.
\]
\end{proof}

\subsection{Exposing the center and routing local tests}

For a fixed center \(x\in\R\), expose its addressed entry in
\emph{every} selector table, including tables outside its route. Write
\begin{equation}
 \label{eq:center-exposure}
 \mathcal F_x=
 \sigma\bigl(S_e(J_{b_e}(x)):e=(P,i),\ i<M\bigr)
 \subseteq\mathcal A.
\end{equation}
An atom \(\xi\) of \(\mathcal F_x\) specifies all these entries and
has positive probability. Their addresses are deterministic once \(x\)
is fixed. On every atom, the remaining selector entries are therefore
independent fair bits, and all terminal entries retain their independent
Bernoulli-\(p\) laws.

\noindent\textit{Remark (notation).} $\xi$ is a fixed positive-probability atom of $\mathcal F_x$, specifying the addressed selector entry at $x$ in every selector table; conditioning on $\xi$ means conditioning on this event. In Lemma~\ref{lem:routing}, Lemma~\ref{lem:independent-tests}, Lemma~\ref{lem:conditional-scales}: $x$ is stable and $\xi$ admits a default choice on the route of $x$; its first default vertex $U$ and height $h$ are fixed on $\xi$.


The exposure determines the entire route of \(x\). If that route uses
a default child, let \(U\) be the first vertex where it does so.
Before conditioning on \(\mathcal F_x\),
\begin{equation}
 \label{eq:no-default}
 \Prob\bigl(\text{the route of \(x\) has no default choice}\bigr)
 =(1-2^{1-M})^d.
\end{equation}
Indeed, reveal the entries along the route one vertex at a time. At the
vertex reached, its \(M-1\) selector entries are independent fair bits,
since the preceding choices use only ancestor tables. The conditional
probability of a default is \(2^{-(M-1)}\). Multiplying the complementary
conditional probabilities over the \(d\) decisions proves the formula.

\noindent\textit{Remark (notation).} $U$ is the $\mathcal F_x$-determined first internal vertex at which the route of $x$ chooses the default child, defined only when that route has a default choice; on any atom $\xi$ admitting $U$, both $U$ and its height $h\in\{1,\ldots,d\}$ are fixed.


Now fix a stable center \(x\) and an atom \(\xi\) on which \(U\)
exists. Under this conditioning, \(U\) and its height \(h\) are fixed.
Put
\[
 r=r_h,\qquad e_i=(U,i)\quad(1\le i<M),\qquad
 y_n(t)=x+tq^n\quad(t\in[1,2]).
\]
All \(M-1\) selector values at \(x\) on these edges are zero. The
windows \(W_{e_i}\) will provide \((M-1)r\) possible hits.

\noindent\textit{Remark (notation).} $e_i=(U,i)$, $1\le i<M$, denotes an outgoing nondefault edge of the first default vertex $U$ of the center $x$; these edges are defined when $U$ exists and are fixed on an atom $\xi$ of $\mathcal F_x$ admitting $U$.


For \(i<M\), let \(L_i(z)\) be the leaf reached by routing \(z\)
starting at the child \(U_i\); if this child is a leaf, take
\(L_i(z)=U_i\). Define the local test
\begin{equation}
 \label{eq:local-test}
 Q_i(z)=S_{e_i}(J_{b_{e_i}}(z))\,
 T_{L_i(z)}(J_{b(L_i(z))}(z)).
\end{equation}
It consults only the selector on \(e_i\) and the tables in its child
subtree. The next lemma makes a successful local test a hit in \(B\).

\noindent\textit{Remark (notation).} For $z\in\R$ and $1\le i<M$, with $U$ fixed by an atom $\xi$ admitting a first default, $L_i(z)$ is the leaf obtained by routing $z$ from $U_i$ using the selectors in that child subtree, with $L_i(z)=U_i$ when $U_i$ is a leaf. For a stable center $x$ and a fixed atom $\xi$ admitting a first default vertex $U$, $Q_i$ is the table-dependent binary function on $\R$ defined by \eqref{eq:local-test}, indexed by $1\le i<M$, with $e_i=(U,i)$ and routing for $L_i(z)$ starting at $U_i$.


\begin{lemma}[Preservation of routing]
 \label{lem:routing}
Fix a stable center \(x\) and an atom \(\xi\) of \(\mathcal F_x\)
on which the first default vertex \(U\) exists. For every outcome in
\(\xi\), every \(i<M\), every \(n\in W_{e_i}\), and every
\(t\in[1,2]\), the route of \(y_n(t)\) reaches \(U\) and rejects its
children \(1,\ldots,i-1\). Consequently,
\begin{equation}
 \label{eq:local-hit}
 Q_i(y_n(t))=1\quad\Longrightarrow\quad y_n(t)\in B.
\end{equation}
\end{lemma}

\begin{proof}
At each strict ancestor \(P\) of \(U\), let \(a\) be the child index
chosen by \(x\). Since \(U\) is its first default vertex, \(a<M\).
The decision at \(P\) reads zeros at indices less than \(a\) and one
at \(a\). All these edges precede \(e_i\) in preorder, so their
selector values at \(y_n(t)\) equal those at \(x\) by
\eqref{eq:stable-keys}. Applying this observation successively from the
root shows that \(y_n(t)\) reaches \(U\). At \(U\), every edge
\((U,j)\) with \(j<i\) also precedes \(e_i\); its selector at
\(y_n(t)\) is therefore the exposed zero at \(x\).

If \(Q_i(y_n(t))=1\), its selector on \(e_i\) is one, so the actual
route takes child \(U_i\). Its continuation reaches the leaf
\(L_i(y_n(t))\), whose addressed terminal entry is one. This is exactly
the membership condition \eqref{eq:random-set}.
\end{proof}

\subsection{The exact failure probability at a fixed scale}

We now show that these tests are numerous enough to make a miss
unlikely. Their routes may share selector entries, so we first fix all
selectors and check that the resulting terminal addresses are distinct.
This order of conditioning keeps the dependence of the leaves on the
selectors explicit.

\begin{lemma}[Independent tests]
 \label{lem:independent-tests}
Fix a stable center \(x\) and an atom \(\xi\) of \(\mathcal F_x\)
on which the first default vertex \(U\) exists. Let \(h\) be its
height and \(r=r_h\), and define \(Q_i\) by \eqref{eq:local-test}.
For every fixed \(t\in[1,2]\),
\begin{equation}
 \label{eq:test-failure}
 \Prob\left(
 Q_i(y_n(t))=0\text{ for all }i<M,\ n\in W_{e_i}
 \,\middle|\,\xi\right)
 =(1-p/2)^{(M-1)r}.
\end{equation}
\end{lemma}

\begin{proof}
On the atom \(\xi\), the index set
\[
 \mathcal I=\{(i,n):1\le i<M,\ n\in W_{e_i}\}
\]
is fixed and has cardinality \((M-1)r\). For each pair put
\(A_{i,n}=S_{e_i}(J_{b_{e_i}}(y_n(t)))\).
For fixed \(i\), Lemma~\ref{lem:separation} gives pairwise distinct
keys for \(x\) and the points \(y_n(t)\), \(n\in W_{e_i}\), at
resolution \(b_{e_i}\). Thus the \(A_{i,n}\) use distinct entries of
their table and avoid its entry exposed at \(x\). Different \(i\)'s
use different tables. Conditional on \(\xi\), all \(A_{i,n}\) are
therefore independent fair bits.

\noindent\textit{Remark (notation).} $\mathcal I=\{(i,n):1\le i<M,\ n\in W_{e_i}\}$, with $e_i=(U,i)$ for the first-default vertex $U$, is fixed on each conditioning atom $\xi$ admitting a first default and has cardinality $(M-1)r_h$, where $h$ is the height of $U$. $A_{i,n}=S_{e_i}(J_{b_{e_i}}(y_n(t)))$ is defined for $1\le i<M$ and $n\in W_{e_i}$. In Lemma~\ref{lem:independent-tests}: For a stable center $x$, an atom $\xi$ of $\mathcal F_x$ admitting a first default, and a scale $t\in[1,2]$ fixed under this conditioning, the index set $\mathcal I=\{(i,n):1\le i<M,\ n\in W_{e_i}\}$ is fixed and has cardinality $(M-1)r_h$; its entries $A_{i,n}$ consult pairwise distinct selector entries, none exposed at $x$.


Next condition on \(\mathcal A\), the sigma-field of all selectors.
For each pair, its local leaf \(L_{i,n}=L_i(y_n(t))\) and terminal
address
\[
 \kappa_{i,n}
 =\bigl(L_{i,n},J_{b(L_{i,n})}(y_n(t))\bigr)
\]
are fixed. These addresses are distinct. If two child indices differ,
the leaves lie in disjoint child subtrees of \(U\), so their tables
differ. If the child indices agree but the leaves differ, the tables
again differ. Finally, suppose two different indices \(n,n'\) below
the same child select the same leaf \(L\). Its incoming edge is either
\(e_i\) itself or a later edge in that child subtree, so
\(b(L)\ge b_{e_i}\). The two points have different keys at resolution
\(b_{e_i}\), and nesting preserves this distinction at resolution
\(b(L)\). Their terminal addresses are distinct in this case as well.

\noindent\textit{Remark (notation).} For $1\le i<M$ and $n\in W_{e_i}$, $L_{i,n}=L_i(y_n(t))$ is the selector-dependent leaf reached by routing $y_n(t)$ from $U_i$. In Lemma~\ref{lem:independent-tests}: With $t$ fixed on $\xi$, conditioning further on $\mathcal A$ fixes $L_{i,n}$. For $1\le i<M$ and $n\in W_{e_i}$, $\kappa_{i,n}=(L_i(y_n(t)),J_{b(L_i(y_n(t)))}(y_n(t)))$ is a selector-dependent terminal address. In Lemma~\ref{lem:independent-tests}: With $t$ fixed on $\xi$, conditioning further on $\mathcal A$ fixes $\kappa_{i,n}$.


Consequently, for each realization of the selectors, the terminal
entries at all addresses \(\kappa_{i,n}\) are independent
Bernoulli-\(p\) variables. The pairs with \(A_{i,n}=0\) fail already;
each remaining pair fails precisely when its terminal bit is zero.
Hence, on \(\xi\),
\[
 \Prob\bigl(Q_i(y_n(t))=0\text{ for all }(i,n)\in\mathcal I
       \mid\mathcal A\bigr)
 =(1-p)^{\sum_{(i,n)\in\mathcal I}A_{i,n}}.
\]
Since \(\mathcal F_x\subseteq\mathcal A\), averaging over the
selectors conditional on \(\xi\) gives
\begin{align*}
 &\Prob\bigl(Q_i(y_n(t))=0\text{ for all }(i,n)\in\mathcal I
       \mid\xi\bigr)\\
 &\qquad=\E\left[(1-p)^{\sum_{(i,n)\in\mathcal I}A_{i,n}}
       \,\middle|\,\xi\right]
 =\prod_{(i,n)\in\mathcal I}\left(\frac12+\frac{1-p}{2}\right)
 =(1-p/2)^{(M-1)r}.
\end{align*}
\end{proof}

The estimate holds for every center-exposure atom admitting a default
and every scale fixed on that atom. Section~\ref{sec:scales} will use
the child-subtree endpoints \(b_{e_i}^*\) to select finitely many such
scales representing all \(t\in[1,2]\). Together with
\eqref{eq:no-default} and Lemma~\ref{lem:density}, this will control
misses while keeping the expected density small.
